
This work establishes that architectural compatibility is a prerequisite for parameter-efficient fine-tuning, not something PEFT methods can create. Zero-shot evaluation provides a principled gate for validating compatibility before compute investment: if a pretrained model fails a task at zero-shot (achieving below-random baseline accuracy), architectural constraints prevent successful fine-tuning regardless of method or hyperparameters.

We demonstrated this principle through systematic evaluation of Mamba-130M, a causal state-space model, on GLUE classification tasks. The model achieves 81\% accuracy on sentiment classification (31 percentage points above random, p < 0.001), demonstrating genuine language understanding when task structure aligns with architectural capabilities. Yet the same checkpoint achieves only 38\% on paraphrase detection (12 percentage points below random, p = 0.003)---a failure indicating that causal processing cannot perform symmetric comparison required for this task. This binary pattern validates our hypothesis: architectural mechanisms determine task feasibility independent of learned parameters.

Our findings challenge assumptions underlying PEFT transfer across architectures. Methods developed for transformers do not automatically transfer to sub-quadratic models. Transformers' bidirectional self-attention enables flexible information flow compatible with diverse tasks, allowing PEFT research to treat the base model as a black box. Causal state-space models and other sub-quadratic architectures introduce constraints that must be validated per-task. Zero-shot evaluation provides this validation with minimal cost---our experiments required under 5 minutes per task, preventing wasted hours of failed fine-tuning.

The zero-shot architectural compatibility gate shifts PEFT workflows from ''fine-tune and hope'' to ''validate then fine-tune.'' Practitioners can test base models on target tasks without gradient updates, interpreting below-random accuracy as a stop signal. This protocol prevents the frustrating cycle of failed hyperparameter searches when the fundamental issue is architectural incompatibility.

As the field explores diverse efficient architectures---linear attention, RWKV, hybrid state-space-transformer models---understanding compatibility boundaries becomes essential. Just as recognizing that architectural alignment enables PEFT success on transformers unlocked efficient adaptation, validating that alignment via zero-shot evaluation prevents failures on misaligned architectures. The future of parameter-efficient fine-tuning lies not just in developing better adaptation methods, but in architectural awareness: understanding which tasks a model can perform before investing in teaching it to perform them well.

Our work opens several research directions. Testing bidirectional state-space models would validate whether dual-direction mechanisms overcome the limitations we identified while preserving efficiency benefits. Developing comprehensive task-architecture compatibility taxonomy would guide practitioners in selecting appropriate base models before experimentation. Exploring architectural augmentation---adding minimal bidirectional layers to causal state-space models---might enable adaptation on previously incompatible tasks while maintaining sub-quadratic complexity.

The broader lesson transcends specific architectures: capabilities must exist before they can be specialized. Fine-tuning optimizes within the constraints of base model design. Zero-shot evaluation exposes those constraints early, guiding research toward viable model-task pairings and preventing wasted effort on fundamentally impossible combinations. In an era of diverse architectural innovation, this principle will only grow more important.
